\chapter{One connection}
\label{ch:connection}

Following \cref{sec:needs}, we take the massive pair and $b$ to be components of one connection, in the minimal orthogonal embedding. The
smallest case is an $\so3$ Yang--Mills field $b^A_\mu$, $A=1,2,3$, with a Higgs vector $n^A$:
\begin{equation}
\label{eq:Lcomp}
L=-\tfrac14F^A_{\mu\nu}F^{A\mu\nu}+\half D_\mu n\cdot D^\mu n-\tfrac\lambda4\,(n\cdot n-v^2)^2,\qquad D_\mu n=\partial_\mu n+g\,b_\mu\times n .
\end{equation}
In the vacuum $n=(0,0,v)$ the component $b^3$ is massless, $b^1$ and $b^2$ have the mass $M=gv$, and the length of $n$ has one massive mode
$h$ of mass $M_h=\sqrt{2\lambda}\,v$. This is the gauge--scalar sector of the model of Georgi and Glashow \cite{GeorgiGlashowPRL72}, an $\so3$
model with an adjoint Higgs field, with the mechanism of \cite{EnglertBrout64,Higgs64}. We write it in the frame variables of Part~I and test its amplitudes; where it lives in the framework is the
subject of \cref{ch:hyp}. Units are $\hbar=c=1$.

\section{The frame form}

\begin{theorem}[Three frame vectors]
\label{thm:frame}
For each component let $X^A=(\bB^A;\bE^A)$ be the frame vector of its field strength, $E_i=F_{0i}$ and $B_i=\half\epsilon_{ijk}F_{jk}$. Then:
\begin{enumerate}[label=(\roman*),leftmargin=2.2em,itemsep=1pt]
\item $L_{\rm YM}=-\half\sum_A\eta(X^A,X^A)$, the sum of three Maxwell Lagrangians of \PI{Prop.~3.20};
\item the Bianchi identities are $D_t\bB^A=\curl_D\bE^A$ and $\dive_D\bB^A=0$: there is no magnetic current;
\item the Yang--Mills equations, in electric and magnetic components, are
\[
\dive_D\bE^A=\rho^A,\qquad D_t\bE^A=-\curl_D\bB^A-\bj^A,\qquad (\rho,\bj)^A=\bigl(-g\,(n\times D_tn)^A,\ g\,(n\times D_in)^A\bigr):
\]
the system of \PI{Prop.~3.10} with covariant derivatives, which carry the non-abelian self-interaction, and with a current made from the
Higgs vector;
\item with the Higgs vector held at its vacuum value $n=(0,0,v)$ and $b$ arbitrary, $g\,(n\times D_\mu n)^A=g^2v^2\,(b^1_\mu,b^2_\mu,0)$, so the
broken components of the source of (iii) satisfy $J^\mu_H=(\rho,\bj)^\mu=-M^2b^\mu_\perp$: a relation of London's form, exact in this gauge.
With the radial mode in unitary gauge, $n=(0,0,v+h)$, $M^2$ becomes $g^2(v+h)^2$.
\end{enumerate}
\checks{co-1, co-2, co-3, co-4}
\end{theorem}

\begin{proposition}[What the completion adds to the working model]
\label{prop:whatadds}
With $W=(b^1+ib^2)/\sqrt2$, $a=b^3$, $f=da$ and $D=\partial+iga$,
\[
-\tfrac14F^A\!\cdot\!F^A=-\tfrac14f^2-\half|D\wedge W|^2+i\,g\,f^{\mu\nu}\bar W_\mu W_\nu
-\tfrac{g^2}2\bigl[(\bar W\!\cdot\!W)^2-(W\!\cdot\!W)(\bar W\!\cdot\!\bar W)\bigr];
\]
for the conjugate field $(b^1-ib^2)/\sqrt2$ the signs of $g$ in $D$ and in the moment term are both reversed. These are minimal coupling of
$W$ to $a$, the magnetic-moment term with coefficient exactly one, that is, gyromagnetic ratio two, and a quartic self-interaction. The
working model of \cref{ch:bcoupling} contains the first of the three terms only.
\checks{co-5}
\end{proposition}

At linear order the massive components of the connection and the massive two-form of \cref{ch:mass} describe the same free propagating
modes in dual variables. With both fields canonically normalised the two-form is the dual curvature of the vector divided by its mass,
$X^a=G\,F^a/M$ up to orientation; in \cref{prop:london} the potential is normalised by $X=G(dC)$ instead. \Cref{prop:whatadds} shows how the
interactions of the two descriptions differ. The identification concerns the free spin-one description; it does not by itself match the
numerical masses. The pole mass of \cref{ch:mass,ch:bcoupling,ch:quantum} equals that of the completion only if $\kappa=g\sqrt Z\,v$, which
with the assignment $v=\kappa$ of \cref{ch:hyp} requires $g\sqrt Z=1$, a choice that \cref{sec:coupled} leaves free.

\section{Amplitudes}

\begin{theorem}[High-energy behaviour at fixed angle]
\label{thm:bounded}
Write $p=P/M$ for the centre-of-mass momentum in units of $M=gv$, and $r=M_h^2/M^2$.
\begin{enumerate}[label=(\roman*),leftmargin=2.2em,itemsep=1pt]
\item The amplitude for a massive quantum scattering off the massless component, computed from \cref{eq:Lcomp} with the vertices extracted
from the Lagrangian, has the cubic and quartic vertices of the vector with unit moment term (\cref{prop:whatadds}), so its growth cancels
at each fixed angle $0<\vartheta<\pi$, with the angle-dependent limit and the backward exception of \cref{prop:moment}; as a numerical
cross-check, its magnitude equals that of the vector amplitude of \cref{prop:moment}, with relative difference below $10^{-8}$, on twelve
random configurations.
\item For $W^+W^-\to W^+W^-$ with all four states longitudinal, at $90^\circ$, away from the pole of the radial exchange,
\begin{equation}
\label{eq:LLLL}
\frac{\mathcal M}{g^2}=2p^2+3-\frac{4(2p^2+1)^2}{4(p^2+1)-r}+\frac{4p^4}{2p^2+r} .
\end{equation}
In the gauge part the contact term is $-(p^2+1)(3p^2+1)$, the neutral exchange in the $s$ channel vanishes at this angle, and that in the
$t$ channel is $3p^4+6p^2+4$; their sum $2p^2+3$ is the amplitude with a Higgs vector of fixed length, $\mathcal M=s/(2v^2)+g^2$, which grows
linearly in $s$. The last two terms of \cref{eq:LLLL} are the exchange of the radial mode in the $s$ and $t$ channels; with them $\mathcal M$
tends to $g^2(3-2r)=3g^2-2M_h^2/v^2$ as $p\to\infty$ for every fixed $r$. At $r=4$, $\mathcal M/g^2=-5+7/p^2-16/p^4+O(p^{-6})$, so its
magnitude approaches the limit from below: the positive powers of the momentum cancel.
\item At $p=100$ and $r=4$ the four contributions are $-300040001$, $+300060004$, $-40004.0001$ and $+19996.0008$, in units of $g^2$, and
their sum is $-4.9993\,g^2$: terms of order $10^8$ cancel to leave the limit of (ii).
\item Over all $81$ helicity combinations at $90^\circ$ with $M_h=2M$ the largest magnitude at $P=300M$ is $5.000\,g^2$ to the displayed
precision; among the channels retained by the numerical threshold the largest exponent fitted between $P=100M$ and $300M$ is about
$10^{-4}$.
\end{enumerate}
\checks{co-6, co-7, co-8, co-9, co-10, co-12, co-13; (ii) the closed form against the vertex extraction at $24$ points and its limit at five
values of $(g,v,\lambda)$; (iv) numerically, at $90^\circ$}
\end{theorem}

Two conditions follow with the criterion $|\mathcal M|<8\pi$. With a Higgs vector of fixed length the leading term $s/(2v^2)$ reaches $8\pi$
at $\sqrt s\approx\sqrt{16\pi}\,v\approx7.1\,v$, of the same order as the working model; within the completion studied, then, the length of the Higgs vector must be a field. With the
radial mode, the $90^\circ$ limit $|3g^2-2M_h^2/v^2|<8\pi$ requires $M_h<\sqrt{(8\pi+3g^2)/2}\;v$, which is $3.75\,v$ at $g=1$ and
$\sqrt{4\pi}\,v\approx3.54\,v$ as $g\to0$, and, when $3g^2>8\pi$, also $M_h^2>(3g^2-8\pi)v^2/2$. These are fixed-angle estimates with our
perturbativity criterion. Lee, Quigg and Thacker derive a partial-wave bound in the electroweak theory \cite{LeeQuiggThacker77}; the
analogous bound here would need normalised partial waves in the coupled channels and a treatment of the massless exchange in the forward
direction. The theory is renormalisable \cite{tHooft71}.

A Higgs vector that points radially, $n=\mathbf x/r$ with $b=0$, gives the field strength of the unbroken charge, the tensor of 't~Hooft
\cite{tHooft74,Polyakov74}, a flux of magnitude $4\pi/g$ through every sphere; with the frame variable $\bB=-\Btext$ and the outward normal
the flux of $\bB$ is $-4\pi/g$. For the global group $SU(2)$, which admits a doublet of charge $g/2$, this is one Dirac unit \cite{Dirac31}, and a probe of charge $g$ measures two.
The configuration is a test of the flux, not a monopole solution: with $b=0$ and fixed length $|\partial_in|^2=2/r^2$, so the gradient
energy in a ball grows linearly with its radius; a smooth monopole of finite energy needs a compensating gauge field and a profile at its
core.
\checks{co-11}

\section{What the completion changes}

Three things change. The mass of the pair is $M=gv$, with $v$ the length of the Higgs vector and its kinetic term normalised as in
\cref{eq:Lcomp}; with a coefficient $Z$ in front of that term, $M=g\sqrt Z\,v$ (\cref{prop:coupled}). Here $v$ is a parameter. Under the
assignment (A3) of \cref{ch:hyp} it is identified with the rate, $v=\kappa$, and the clock angle $\theta$ of \cref{thm:clock} then
corresponds to the angle of the rotation generated by the generator of mass (\cref{thm:massgen}). The correspondence is kinematic: the
action \cref{eq:Lcomp} contains no $\theta$, and a dynamical identification would need an additional coupling or field equation.

Next, each component of the curvature is a frame vector, and its Lagrangian and field equations are those of \PI{Props.~3.10 and~3.20}
with covariant derivatives (\cref{thm:frame}). \PI{Section~8.6} leaves open the link between the field-strength reading and the internal
connection; here the link takes this form for the algebra of this chapter. It is an identification, not a derivation: the Yang--Mills term,
the kinetic term and the potential of the Higgs vector are assumed.

Finally, the magnetic current disappears. In \cref{ch:mass} the field strength of the two-form was dual to the magnetic current, while in
the completed theory the Bianchi identity removes it, and the current that remains is the electric one, made from the Higgs vector.
